% Delta focal 2026-09-26. Los cuerpos 11/12 del catalogo permanecen inalterados.
\section{Longitud estructural, acci\'on y respuesta gravitatoria}
\label{sec:vi-catalogo-dependencia-l-g}

El cat\'alogo re\'une 324 rutas y 855 registros metrol\'ogicos. Su
interpretaci\'on sigue la construcci\'on del mismo estado HMT--MD, en el orden
causal
\[
 \mathrm{APP}\longrightarrow\mathrm{TRIT}\longrightarrow\mathrm{TPK}
 \longrightarrow X_{\mathrm{enr}}
 \longrightarrow\mathcal C^{\mathrm{disc}}_{\mathrm{joint}}
 \longrightarrow L_{\alpha}
 \longrightarrow(H\Lambda,R\Lambda)
 \longrightarrow G .
 \tag{A.1}\label{eq:vi-cat-cadena-lg}
\]
Las coordenadas \(\pi\), \(\varphi\), \(e\) y
\(\alpha_{\mathrm{HMT}}\) son salidas correlacionadas de la imagen non\'adica
tipada anterior; no se introducen como valores convencionales para seleccionar
esta realizaci\'on. Para \(N=12\cdot4\cdot120=5760\), la longitud marcada es
\[
 L_{\alpha}=\alpha_{\mathrm{HMT}}^{16}
 \frac{N-1}{4N}L_*
 =\frac{5759}{23040}\alpha_{\mathrm{HMT}}^{16}L_* .
 \tag{A.2}\label{eq:vi-cat-longitud-alpha}
\]

Abreviamos en esta secci\'on \(\hbar=\hbar_{\mathrm{ret}}\) y
\(c=c_{\mathrm{int}}\). Sobre la fibra positiva no nula, la energ\'ia
\(H\), la longitud de acci\'on reducida \(\Lambda\) y la lectura radial
conjugada \(R\) satisfacen las dos relaciones operatorias siguientes.
Definimos adem\'as \(M:=H/c^2\) y declaramos la identificaci\'on f\'isica
radial \(R:=GM/c^2\). Al compararla, sobre la fibra invertible, con la
relaci\'on operatoria \(R=L_{\alpha}^{2}H/(\hbar c)\), se obtiene
\[
 H\Lambda=\hbar c\,I,\qquad
 R\Lambda=L_{\alpha}^{2}I,\qquad
 R=\frac{L_{\alpha}^{2}}{\hbar c}H,\qquad
 \boxed{G=\frac{c^{3}L_{\alpha}^{2}}{\hbar}} .
 \tag{A.3}\label{eq:vi-cat-realizacion-lg}
\]
El reconocimiento metrol\'ogico de \(G\) es posterior y no interviene en la
elecci\'on del estado, del lector ni de sus coeficientes.

\subsection{Registro M10 y conservaci\'on de la ley de masas}
\label{sec:vi-catalogo-m10-activo}

El c\'odigo M10 del repertorio documental remite a la
realizaci\'on tipada de \eqref{eq:vi-cat-realizacion-lg}. M10 no postula una
part\'icula gravitatoria ni convierte en salida HMT el registro externo que el
inventario conserva bajo la clase
\texttt{GRAVITATIONAL\_OR\_HYPOTHETICAL}: dicho registro permanece como dato
metrol\'ogico externo, separado de la construcci\'on
\(L_{\alpha}\mapsto G\).

La escala energ\'etica inducida por la misma longitud es
\[
 Q_L:=\frac{\hbar c}{L_{\alpha}}.
 \tag{A.4}\label{eq:vi-cat-escala-q}
\]
Si una coordenatización equivalente sustituye \(Q_L\) por \(Q_L'\), el coeficiente
electr\'onico cotransforma seg\'un
\[
 b_e'=b_e\frac{Q_L}{Q_L'},\qquad Q_L'b_e'=Q_Lb_e .
 \tag{A.5}\label{eq:vi-cat-cotransformacion-be}
\]
Las definiciones de la energ\'ia del reloj y del factor electr\'onico dan
\[
 E_{\mathrm{clk}}=Q_L,\qquad
 m_{\mathrm{clk}}=\frac{E_{\mathrm{clk}}}{c^2},\qquad
 b_e=\frac{E_e^{(0)}}{E_{\mathrm{clk}}},\qquad
 m_{\mathrm{clk}}b_e=\frac{E_e^{(0)}}{c^2}.
\]
La cancelaci\'on exacta de \(Q_L\) demuestra que la realizaci\'on de \(G\) no
deriva de nuevo las masas, no selecciona rutas y no modifica las tablas ya
construidas. La especializaci\'on \(\Xi_\Gamma=1\) del teorema de
compatibilidad equivale a
\(\bar\lambda_\Gamma=R_{108}=L_{\alpha}\): identifica una coincidencia de
longitudes y no afirma ni crea la existencia de una especie de masa de Planck
en el cat\'alogo.

La d\'ecada \(10^{-34}\) pertenece exclusivamente a la rama absoluta
acci\'on--reloj--masa,
\(\hbar=\eta_{\mathrm{ret}}10^{-34}\mathcal U_S\); no se reinserta en el
exponente relativo, en el car\'acter de ruta ni como correcci\'on por especie.

\subsection{Canales constitutivos y momento orientado de Catal\'an}
\label{sec:vi-catalogo-propietario-constitutivo}

Sean \(0<q_+,q_-<1\) los canales etiquetados de transporte y \(P_+,P_-\)
sus proyectores ortogonales de hoja. Para \(s\in(0,1)\), definimos
\[
 \mathcal C_2(s)=\sum_{k\geq0}
 \frac{(-1)^k s^{2k+1}}{(2k+1)^2},\qquad
 \mathcal D=s\frac{\mathrm d}{\mathrm ds},\qquad
 f(s)=\frac{1-s^3}{1-s^4}.
 \tag{A.6}\label{eq:vi-cat-constitutivas}
\]
La convergencia local uniforme de la serie y de sus derivadas permite aplicar
dos veces \(\mathcal D\). La serie geom\'etrica resultante da
\[
 \mathcal D^2\mathcal C_2(s)=\frac{s}{1+s^2},\qquad
 f(s)=\frac{1+s+s^2}{s(1+s)}\mathcal D^2\mathcal C_2(s).
\]
Las identidades \(90=3\cdot30\) y \(120=4\cdot30\) determinan
\[
 s_{\pm}=q_{\pm}^{30},\qquad
 r_{\pm}=f(s_{\pm})
 =\frac{1-q_{\pm}^{90}}{1-q_{\pm}^{120}}.
\]
Por tanto, \(V=r_+P_++r_-P_-\) re\'une las respuestas conservando sus
etiquetas. La serie de \(\mathcal C_2\) converge absolutamente en \(s=1\);
su valor es \(G_{\mathrm{Cat}}\), distinguido del coeficiente gravitatorio
\(G\).

La funci\'on completa se recupera de \(f\) mediante
\[
 \mathcal C_2(s)=\int_0^s\log\!\left(\frac{s}{t}\right)
 \frac{1+t}{1+t+t^2}f(t)\,\mathrm dt.
\]
En efecto, \((1+t)f(t)/(1+t+t^2)=1/(1+t^2)\); dos aplicaciones de
\(\mathcal D\) restituyen \(s/(1+s^2)\). Las soluciones de la ecuaci\'on
homog\'enea son \(a\log s+b\); la condici\'on
\(\mathcal C_2(s)\to0\) cuando \(s\downarrow0\) fija ambos coeficientes a
cero. Quedan expl\'icitos la respuesta, su restituci\'on regular y el
significado de la coordenada de Catal\'an usada en las fichas.
